\SourceLeaf{14}
\begin{center}
{\Large\textsc{UNTERSUCHUNGEN ÜBER HÖHERE ARITHMETIK.}}

\BookRule{\linewidth}

\textsc{ERSTER ABSCHNITT}

{\scriptsize\textsc{ÜBER}}

\textsc{DIE KONGRUENZ VON ZAHLEN IM ALLGEMEINEN.}

\BookRule{0.10\linewidth}
\end{center}

\Needspace{5\baselineskip}
\begin{center}
\emph{Kongruente Zahlen, Moduln, Reste und Nichtreste.}
\end{center}

\Article{1}

Wenn die Zahl $a$ die Differenz der Zahlen $b, c$ teilt, heißen $b$ und $c$ \emph{kongruent nach $a$}, andernfalls \emph{inkongruent}: $a$ selbst nennen wir den \emph{Modul}. Im ersten Fall heißt jede der Zahlen $b, c$ ein \emph{Rest} der anderen, im zweiten Fall ein \emph{Nichtrest}.

Diese Begriffe gelten für alle ganzen Zahlen, sowohl positive als auch negative\HistoricalNote{*)}{Der Modul ist immer \emph{absolut}, also ohne Vorzeichen, zu nehmen.}, sollen aber nicht auf Brüche ausgedehnt werden. \emph{Zum Beispiel} sind $-9$ und $+16$ nach dem Modul $5$ kongruent; $-7$ ist ein Rest von $+15$ nach dem Modul $11$, nach dem Modul $3$ dagegen ein Nichtrest. Da außerdem jede Zahl die Null teilt, ist jede Zahl nach jedem Modul als zu sich selbst kongruent anzusehen.

\Article{2}

Alle Reste einer gegebenen Zahl $a$ nach dem Modul $m$ sind in der Formel $a+km$ enthalten, wobei $k$ eine unbestimmte ganze Zahl bezeichnet. Die einfacheren Sätze, die wir später angeben werden, lassen sich daraus ohne Mühe beweisen; ihre Richtigkeit kann aber jeder ebenso leicht unmittelbar erkennen.

\SourceLeaf{15}
Von nun an werden wir die Kongruenz von Zahlen durch das Zeichen $\equiv$ bezeichnen und, wo es nötig ist, den Modul in Klammern hinzufügen, $-16\equiv 9\ (\mathrm{mod.}\ 5)$. $-7\equiv 15\ (\mathrm{mod.}\ 11)$\HistoricalNote{*)}{Wir haben dieses Zeichen wegen der engen Analogie zwischen Gleichheit und Kongruenz gewählt. Aus demselben Grund hat der ausgezeichnete Le Gendre in der Abhandlung, die wir weiter unten öfter anführen werden, für die Kongruenz das Gleichheitszeichen selbst beibehalten; wir haben gezögert, ihm darin zu folgen, um Mehrdeutigkeiten zu vermeiden.}.

\Article{3}

\textsc{Satz.} \emph{Sind $m$ aufeinanderfolgende ganze Zahlen}
\[
a,\ a+1,\ a+2\ \GaussFourDots\ a+m-1
\]
\emph{und eine weitere ganze Zahl $A$ gegeben, so ist genau eine der ersteren nach dem Modul $m$ zu dieser kongruent.}

Ist nämlich $\frac{a-A}{m}$ ganzzahlig, so ist $a\equiv A$; ist es ein Bruch, so sei die nächstgrößere ganze Zahl (oder, wenn er negativ ist, die nächst\emph{kleinere}, wenn man vom Vorzeichen absieht) $=k$. Dann liegt $A+km$ zwischen $a$ und $a+m$ und ist also die gesuchte Zahl. Alle Quotienten $\frac{a-A}{m}$, $\frac{a+1-A}{m}$, $\frac{a+2-A}{m}$ usw. liegen offensichtlich zwischen $k-1$ und $k+1$; daher kann nicht mehr als einer von ihnen ganzzahlig sein.

\Needspace{5\baselineskip}
\begin{center}\emph{Kleinste Reste.}\end{center}
\Article{4}

Jede Zahl hat also sowohl in der Reihe $0, 1, 2,\ldots m-1$ als auch in der Reihe $0, -1, -2,\ \GaussFourDots\ -(m-1)$ einen Rest; diese nennen wir \emph{kleinste Reste}. Ist nicht $0$ ein Rest, so gibt es offensichtlich immer zwei, einen \emph{positiven} und einen \emph{negativen}. Sind sie dem Betrag nach verschieden, so ist einer $<\frac{m}{2}$; andernfalls sind beide $=\frac{m}{2}$, wenn man vom Vorzeichen absieht. Daher hat jede Zahl einen Rest, der die Hälfte des Moduls nicht übersteigt; dieser heißt \emph{absolut kleinster Rest}.

\emph{Zum Beispiel} hat $-13$ nach dem Modul $5$ den kleinsten positiven Rest $2$, der zugleich der absolut kleinste ist, und den kleinsten negativen Rest $-3$; $+5$ ist nach dem Modul $7$ sein eigener kleinster positiver Rest, während $-2$ der kleinste negative und zugleich der absolut kleinste Rest ist.

\Needspace{5\baselineskip}
\begin{center}\emph{Elementare Sätze über Kongruenzen.}\end{center}
\Article{5}

Nachdem diese Begriffe feststehen, wollen wir die Eigenschaften kongruenter Zahlen zusammenstellen, die sich unmittelbar zeigen.

\SourceLeaf{16}
\emph{Zahlen, die nach einem zusammengesetzten Modul kongruent sind, sind auch nach jedem seiner Teiler kongruent.}

\emph{Sind mehrere Zahlen nach demselben Modul zu derselben Zahl kongruent, so sind sie auch untereinander kongruent} (nach demselben Modul).

Auch im Folgenden ist stillschweigend vorausgesetzt, dass der Modul derselbe ist.

\emph{Kongruente Zahlen haben dieselben kleinsten Reste, inkongruente verschiedene.}

\Article{6}

\emph{Sind beliebig viele ganze Zahlen $A, B, C$ usw. und ebenso viele andere $a, b, c$ usw. gegeben, die ihnen nach irgendeinem Modul kongruent sind,}
\[
A\equiv a,\ B\equiv b\ \text{\itshape etc.,\quad gilt}\quad A+B+C+\text{\itshape etc.}\equiv a+b+c+\text{\itshape etc.}
\]

\emph{Ist $A\equiv a, B\equiv b$, so ist $A-B\equiv a-b$.}

\Article{7}

\emph{Ist $A\equiv a$, so ist auch $kA\equiv ka$.}

Ist $k$ eine positive Zahl, so ist dies nur ein besonderer Fall des Satzes im vorigen Artikel, wenn man dort $A=B=C$ usw., $a=b=c$ usw. setzt. Ist $k$ negativ, so ist $-k$ positiv, also $-kA\equiv -ka$, und daraus folgt $kA\equiv ka$.

\emph{Ist $A\equiv a, B\equiv b$, so ist $AB\equiv ab$.} Denn $AB\equiv Ab\equiv ba$.

\Article{8}

\emph{Sind beliebig viele ganze Zahlen $A, B, C$ usw. und ebenso viele ihnen kongruente Zahlen $a, b, c$ usw. gegeben, $A\equiv a, B\equiv b$ usw., so sind die Produkte beider Gruppen kongruent, $ABC\ \text{\itshape etc.}\equiv abc\ \text{\itshape etc.}$}

Nach dem vorigen Artikel ist $AB\equiv ab$ und aus demselben Grund $ABC\equiv abc$; auf dieselbe Weise können beliebig viele weitere Faktoren hinzukommen.

Nimmt man alle Zahlen $A, B, C$ usw. als gleich an und ebenso die entsprechenden Zahlen $a, b, c$ usw., so erhält man den Satz: \emph{Ist $A\equiv a$ und $k$ eine positive ganze Zahl, so ist $A^{k}\equiv a^{k}$.}

\Article{9}

\emph{Sei $\mathrm{X}$ eine algebraische Funktion der Unbestimmten $x$ von der Form}
\[
Ax^{a}+Bx^{b}+Cx^{c}+\text{\itshape etc.}
\]
\emph{wobei $A, B, C$ usw. beliebige ganze Zahlen und $a, b, c$ usw. nichtnegative ganze Zahlen bezeichnen. Gibt man dann der Unbestimmten $x$ nach irgendeinem Modul kongruente Werte, so sind auch die daraus entstehenden Werte der Funktion $\mathrm{X}$ kongruent.}

\SourceLeaf{17}
Seien $f, g$ kongruente Werte von $x$. Dann ist nach dem vorigen Artikel $f^{a}\equiv g^{a}$ und $Af^{a}\equiv Ag^{a}$, ebenso $Bf^{b}\equiv Bg^{b}$ usw. Daher
\[
Af^{a}+Bf^{b}+Cf^{c}+\text{etc.}\equiv Ag^{a}+Bg^{b}+Cg^{c}+\text{etc.}\quad\text{\itshape Q. E. D.}
\]

Es ist außerdem leicht zu sehen, wie sich dieser Satz auf Funktionen mehrerer Unbestimmter ausdehnen lässt.

\Article{10}

Setzt man also für $x$ alle aufeinanderfolgenden ganzen Zahlen ein und führt die Werte der Funktion $\mathrm{X}$ auf ihre kleinsten Reste zurück, so bilden diese eine Reihe, in der dieselben Glieder nach einem Abstand von $m$ Gliedern wiederkehren (wobei $m$ den Modul bezeichnet); die Reihe besteht also aus einer unendlich oft wiederholten \emph{Periode} von $m$ Gliedern. Sei zum Beispiel $\mathrm{X}=x^{3}-8x+6$ und $m=5$; dann ergeben für $x=0, 1, 2, 3$ usw. die Werte von $\mathrm{X}$ die kleinsten positiven Reste $1, 4, 3, 4, 3, 1, 4$ usw., wobei sich die ersten fünf $1, 4, 3, 4, 3$ unbegrenzt wiederholen. Setzt man die Reihe rückwärts fort, gibt man also $x$ negative Werte, so erscheint dieselbe Periode in umgekehrter Reihenfolge ihrer Glieder: daher können in der ganzen Reihe keine anderen Glieder vorkommen als die dieser Periode.

\Article{11}

In diesem Beispiel kann also $\mathrm{X}$ weder $\equiv 0$ noch $\equiv 2$ (mod. 5) werden, erst recht nicht $=0$ oder $=2$. Folglich lassen sich die Gleichungen $x^{3}-8x+6=0$ und $x^{3}-8x+4=0$ nicht durch ganze Zahlen und somit, wie bekannt, auch nicht durch rationale Zahlen lösen. Allgemein ist klar, dass die Gleichung $\mathrm{X}=0$, wenn $\mathrm{X}$ eine Funktion der Unbekannten $x$ von der Form
\[
x^{n}+Ax^{n-1}+Bx^{n-2}+\text{etc.}+N
\]
mit ganzen Zahlen $A, B, C$ usw. und einer positiven ganzen Zahl $n$ ist (auf diese Form lassen sich bekanntlich alle algebraischen Gleichungen bringen), keine rationale Wurzel hat, wenn sich die Kongruenz $\mathrm{X}\equiv 0$ nach irgendeinem Modul nicht erfüllen lässt. Dieses Kriterium, das sich uns hier von selbst angeboten hat, wird in Abschnitt VIII ausführlicher behandelt. Dieses Beispiel kann gewiss zumindest eine Vorstellung vom Nutzen dieser Untersuchungen geben.

\SourceLeaf{18}
\Needspace{5\baselineskip}
\begin{center}\emph{Einige Anwendungen.}\end{center}
\Article{12}

Auf den Sätzen dieses Kapitels beruhen mehrere Regeln, die man gewöhnlich in der Arithmetik lehrt, zum Beispiel die Regeln, mit denen man prüft, ob eine gegebene Zahl durch 9, 11 oder andere Zahlen teilbar ist. \emph{Nach dem Modul} 9 sind alle Potenzen der Zahl 10 zu eins kongruent: hat also die gegebene Zahl die Form $a+10b+100c+\text{etc.}$, so liefert sie nach dem Modul 9 denselben kleinsten Rest wie $a+b+c+\text{etc.}$ Addiert man also die einzelnen Ziffern einer in Dezimalschreibweise dargestellten Zahl unabhängig von ihrer Stelle, so liefern diese Summe und die gegebene Zahl dieselben kleinsten Reste. Die gegebene Zahl ist daher genau dann durch 9 teilbar, wenn die Summe durch 9 teilbar ist. Dasselbe gilt für den Teiler 3. Da \emph{nach dem Modul} 11 $100\equiv 1$ ist, gilt allgemein $10^{2k}\equiv 1$. $10^{2k+1}\equiv 10\equiv -1$, und eine Zahl von der Form $a+10b+100c+\text{etc.}$ liefert nach dem Modul 11 denselben kleinsten Rest wie $a-b+c$ usw.; daraus folgt sofort die bekannte Regel. Alle ähnlichen Regeln lassen sich aus demselben Grundsatz leicht ableiten.

Die vorigen Ergebnisse erklären ebenso die gewöhnlich empfohlenen Regeln zur Prüfung arithmetischer Rechnungen. Sollen nämlich aus gegebenen Zahlen durch Addition, Subtraktion, Multiplikation oder Potenzieren andere Zahlen abgeleitet werden, so ersetzt man die gegebenen Zahlen durch ihre kleinsten Reste nach einem beliebigen Modul (gewöhnlich 9 oder 11, weil sich die Reste nach diesen Moduln, wie wir eben gezeigt haben, in unserem Dezimalsystem so leicht finden lassen). Die so entstehenden Zahlen müssen zu den aus den gegebenen Zahlen abgeleiteten kongruent sein; trifft dies nicht zu, so hat sich ein Fehler in die Rechnung eingeschlichen.

Da diese und ähnliche Dinge aber hinreichend bekannt sind, wäre es überflüssig, länger dabei zu bleiben.

\par\bigskip\bigskip\BookRule{0.29\linewidth}
